\documentclass[10pt]{extarticle}
\usepackage[a4paper,margin=1.5cm]{geometry}
\usepackage[T1]{fontenc}
\usepackage{amsmath,amssymb,mathtools}
\usepackage{graphicx}
\usepackage[export]{adjustbox}
\usepackage{multicol}
\usepackage{enumitem}
\usepackage{booktabs,array,tabularx}
\setlength{\columnsep}{0.7cm}
\setlength{\parindent}{0pt}
\setlength{\parskip}{3pt}
\pagestyle{empty}
\makeatletter
\newcommand{\qfigdef}[4]{\@namedef{qf@file@#1}{#2}\@namedef{qf@w@#1}{#3}\@namedef{qf@wide@#1}{#4}}
\newcommand{\qf@check}[1]{\@ifundefined{qf@file@#1}{\PackageError{score2}{Unknown figure id #1}{}}{}}
\newcommand{\qf@img}[4][]{\includegraphics[width=#2\dimexpr\@nameuse{qf@w@#4}pt\relax,max width=#3,max height=0.45\textheight,#1]{\@nameuse{qf@file@#4}}}
\newif\ifqf@full
\newif\ifqf@boxed
\newsavebox{\qf@box}
\newcommand{\qf@open}{\global\qf@fulltrue\par\global\qf@boxedtrue
  \setbox\qf@box\vbox\bgroup\hsize\textwidth\linewidth\textwidth}
\newcommand{\qf@emit}[1]{\par\egroup\global\qf@boxedfalse
  \ifdim\dimexpr\ht\qf@box+\dp\qf@box\relax<#1\textheight\noindent\usebox{\qf@box}\par\else\unvbox\qf@box\fi}
\newcommand{\qfullsplit}{\ifqf@boxed\qf@emit{0.85}\fi}
\newcommand{\qf@close}{\ifqf@boxed\qf@emit{0.45}\fi\global\qf@fullfalse}
\newcommand{\qfullbegin}{\end{multicols}\qf@open}
\newcommand{\qfullend}{\qf@close\begin{multicols}{2}}
\newcommand{\qfullstart}{\qf@open}
\newcommand{\qfulldone}{\qf@close}
\newcommand{\qfig}[2][1]{\qf@check{#2}%
  \ifqf@full\qf@blk{#1}{#2}\else
  \ifnum\@nameuse{qf@wide@#2}=1 \expandafter\@firstoftwo\else\expandafter\@secondoftwo\fi
  {\end{multicols}\begin{center}\qf@img{#1}{\textwidth}{#2}\end{center}\begin{multicols}{2}}%
  {\qf@blk{#1}{#2}}\fi}
\newcommand{\qf@blk}[2]{\par\smallskip{\centering\qf@img{#1}{\linewidth}{#2}\par}\smallskip}
\newcommand{\qfigi}[2][1]{\qf@check{#2}\qf@img[valign=c]{#1}{\linewidth}{#2}}
\makeatother
\newcommand{\papersec}[1]{\par\bigskip{\centering\Large\bfseries PAPER-#1\par}\medskip}
\newcommand{\partsec}[2]{\par\medskip\noindent{\large\bfseries PART-#1: #2}\par\smallskip}
\newcommand{\sect}[1]{\par\smallskip\noindent\textbf{#1}\par}
\newcommand{\qn}[1]{\par\smallskip\noindent\textbf{#1.}~}
\newcommand{\paraq}[2]{\par\smallskip\noindent\textbf{Paragraph for Questions #1 and #2}\par}
\newenvironment{qoptions}{\begin{itemize}[leftmargin=1.8em,itemsep=1pt,topsep=2pt,label={}]}{\end{itemize}}
\newcommand{\opt}[2]{\item[(#1)] #2}
\newcolumntype{L}[1]{>{\raggedright\arraybackslash}p{#1}}
\newcommand{\fwstart}{\end{multicols}\par\noindent\begin{minipage}{\textwidth}\small\setlength{\tabcolsep}{4pt}\renewcommand{\arraystretch}{1.5}}
\newcommand{\fwend}{\end{minipage}\par\medskip\begin{multicols}{2}}
% figure registry (generated by score2_tools.py finalize)
\qfigdef{p1_phy_s1_q01_1}{SRG (Enthusiast) - Test 03 (18 Nov 2020)-img/p1_phy_s1_q01_1.png}{132.0}{0}
\qfigdef{p1_phy_s2_q09_A}{SRG (Enthusiast) - Test 03 (18 Nov 2020)-img/p1_phy_s2_q09_A.png}{69.4}{0}
\qfigdef{p1_phy_s2_q09_B}{SRG (Enthusiast) - Test 03 (18 Nov 2020)-img/p1_phy_s2_q09_B.png}{69.4}{0}
\qfigdef{p1_phy_s2_q09_C}{SRG (Enthusiast) - Test 03 (18 Nov 2020)-img/p1_phy_s2_q09_C.png}{69.4}{0}
\qfigdef{p1_phy_s2_q09_D}{SRG (Enthusiast) - Test 03 (18 Nov 2020)-img/p1_phy_s2_q09_D.png}{69.6}{0}
\qfigdef{p1_phy_s2_q11_1}{SRG (Enthusiast) - Test 03 (18 Nov 2020)-img/p1_phy_s2_q11_1.png}{122.4}{0}
\qfigdef{p1_phy_s3_q14_1}{SRG (Enthusiast) - Test 03 (18 Nov 2020)-img/p1_phy_s3_q14_1.png}{122.8}{0}
\qfigdef{p1_chem_s2_q07_1}{SRG (Enthusiast) - Test 03 (18 Nov 2020)-img/p1_chem_s2_q07_1.png}{193.5}{0}
\qfigdef{p1_chem_s2_q07_A}{SRG (Enthusiast) - Test 03 (18 Nov 2020)-img/p1_chem_s2_q07_A.png}{318.4}{1}
\qfigdef{p1_chem_s2_q09_1}{SRG (Enthusiast) - Test 03 (18 Nov 2020)-img/p1_chem_s2_q09_1.png}{207.3}{0}
\qfigdef{p1_chem_s2_q10_1}{SRG (Enthusiast) - Test 03 (18 Nov 2020)-img/p1_chem_s2_q10_1.png}{185.5}{0}
\qfigdef{p1_chem_s2_q12_1}{SRG (Enthusiast) - Test 03 (18 Nov 2020)-img/p1_chem_s2_q12_1.png}{352.2}{1}
\raggedcolumns\makeatletter\@beginparpenalty=10000\makeatother

\begin{document}
\begin{center}{\Large\bfseries SRG (Enthusiast) --- Test 03 (18 Nov 2020)}\end{center}
\begin{multicols}{2}
\partsec{1}{PHYSICS}
\sect{SECTION-I (i)}
\qn{1} In a meter bridge, null point is at $\ell=33.7$ cm, when the resistance $S$ is shunted by 12 $\Omega$ resistance the null point is found to be shifted by a distance of 18.2 cm. The value of unknown resistance $R$ should be:
\qfig{p1_phy_s1_q01_1}
\begin{qoptions}
\opt{A}{13.5 $\Omega$}
\opt{B}{68.8 $\Omega$}
\opt{C}{3.42 $\Omega$}
\opt{D}{None of these}
\end{qoptions}
\qn{2} In optical bench experiment index error is +1 cm and $-1$ cm, between object needle and lens, lens and image needle respectively. Observed values of $u$ and $\upsilon$ are 9 cm and 17 cm. Focal length of lens is:
\begin{qoptions}
\opt{A}{6.15 cm}
\opt{B}{5.54 cm}
\opt{C}{5.88 cm}
\opt{D}{6.25 cm}
\end{qoptions}
\qn{3} The pitch of a screw gauge is 0.5 mm and there are 100 divisions on it circular scale. The instrument reads +2 divisions when nothing is put in-between its jaws. In measuring the diameter of a wire, there are 8 divisions on the main scale and 83$^{rd}$ division coincides with the reference line. Then the diameter of the wire is:
\begin{qoptions}
\opt{A}{4.05 mm}
\opt{B}{4.405 mm}
\opt{C}{3.05 mm}
\opt{D}{1.25 mm}
\end{qoptions}
\qn{4} Radiation of two photon energies twice and five times the work function of metal are incident successively on the metal surface. The ratio of the maximum velocity of photoelectrons emitted in the two cases will be:
\begin{qoptions}
\opt{A}{1 : 2}
\opt{B}{2 : 1}
\opt{C}{1 : 4}
\opt{D}{4 : 1}
\end{qoptions}
\sect{SECTION-I (ii)}
\qn{5} In a YDSE arrangement, white light is used to illuminate the slits. At point on the screen directly in front of the slits, which of these wavelengths is/are missing?

($\lambda$ = wavelength of light used, $d$ = distance between the slits, $D$ = separation between the slits and the screen, $d\ll D$):
\begin{qoptions}
\opt{A}{$\dfrac{d^2}{D}$}
\opt{B}{$\dfrac{d^2}{3D}$}
\opt{C}{$\dfrac{d^2}{4D}$}
\opt{D}{$\dfrac{d^2}{8D}$}
\end{qoptions}
\qn{6} A thin paper of thickness 0.02 mm having refractive index 1.45 is pasted across one of the slits in a Young's double slit experiment. The paper transmits $\dfrac49$ of light falling on it ($\lambda_{\text{light}}=600$ nm):
\begin{qoptions}
\opt{A}{Amplitude of light wave transmitted through the paper will be $\dfrac23$ times that of incident wave.}
\opt{B}{The ratio of maximum and minimum intensity in the fringe pattern will be 25.}
\opt{C}{The total number of fringe crossing the centre if an identical paper is pasted on the other slit is 15.}
\opt{D}{The ratio of maximum and minimum intensity in the pattern will be 5.}
\end{qoptions}
\qn{7} If radiation of all wavelengths from ultraviolet to infrared is passed through hydrogen gas at room temperature, absorption lines will be observed in the:
\begin{qoptions}
\opt{A}{Lyman series}
\opt{B}{Balmer series}
\opt{C}{Both (A) and (B)}
\opt{D}{Neither (A) nor (B)}
\end{qoptions}
\qn{8} When a nucleus with atomic number $Z$ and mass number $A$ undergoes a radioactive decay process:
\begin{qoptions}
\opt{A}{Both $Z$ and $A$ will decrease, if the process is $\alpha$ decay}
\opt{B}{$Z$ will decrease but $A$ will not change, if the process is $\beta^+$ decay}
\opt{C}{$Z$ will decrease but $A$ will not change, if the process is $\beta^-$ decay}
\opt{D}{$Z$ and $A$ will remain unchanged, if the process is $\gamma$ decay}
\end{qoptions}
\qn{9} A point source causes photoelectric effect from a small metal plate. Which of the following curves may represent the saturation photocurrent as a function of the distance between the source and the metal?
\begin{qoptions}
\opt{A}{\qfigi{p1_phy_s2_q09_A}}
\opt{B}{\qfigi{p1_phy_s2_q09_B}}
\opt{C}{\qfigi{p1_phy_s2_q09_C}}
\opt{D}{\qfigi{p1_phy_s2_q09_D}}
\end{qoptions}
\qn{10} By increasing the intensity of incident light keeping frequency $(\upsilon>\upsilon_0)$ fixed on the surface of metal:
\begin{qoptions}
\opt{A}{Kinetic energy of the photoelectrons increases}
\opt{B}{Number of emitted electrons increases}
\opt{C}{Kinetic energy and number of electrons increases}
\opt{D}{No effect}
\end{qoptions}
\qn{11} A young's double slit experiment is conducted in water $(\mu_1)$ as shown in the figure, and a glass plate of thickness $t$ and refractive index $\mu_2$ is placed in the path of $S_2$. The magnitude of the phase difference at O is:

(Assume that `$\lambda$' is the wavelength of light in air). O is symmetrical w.r.t. $S_1$ and $S_2$.
\qfig{p1_phy_s2_q11_1}
\begin{qoptions}
\opt{A}{$\left|\left(\dfrac{\mu_2}{\mu_1}-1\right)t\right|\dfrac{2\pi}{\lambda}$}
\opt{B}{$\left|\left(\dfrac{\mu_1}{\mu_2}-1\right)t\right|\dfrac{2\pi}{\lambda}$}
\opt{C}{$\left|(\mu_2-\mu_1)t\right|\dfrac{2\pi}{\lambda}$}
\opt{D}{$\left|(\mu_2-1)t\right|\dfrac{2\pi}{\lambda}$}
\end{qoptions}
\qn{12} The photoelectric effect can not be explained by the wave theory of light because:
\begin{qoptions}
\opt{A}{The energy carried by the light waves is not given to a particular electron of the metal, rather it is distributed among all the electrons present on the surface of metal}
\opt{B}{Waves do not have energy}
\opt{C}{Energy of the waves becomes zero as it strikes the metal surface}
\opt{D}{Waves do not have sufficient energy which is required for electron emission}
\end{qoptions}
\sect{SECTION-II}
\qn{13} The shortest wavelength of the Bracket series of a hydrogen - like atom (atomic number = $Z$) is the same as the shortest wavelength of the Balmer series of hydrogen atom. Find the value of $Z$.
\qn{14} A free particle of mass `$m$' is confined to a region of length $L$. The de-Broglie wave associated with the particle is sinusoidal in nature as given in figure. The kinetic energy of the particle is $\dfrac{n^2h^2}{8mL^2}$ then find the value of $n$.
\qfig{p1_phy_s3_q14_1}
\qn{15} One slit of a Young's experiment is covered by a glass plate $(\mu_1=1.4)$ and the other by another glass plate $(\mu_2=1.7)$ of the same thickness. The point of central maxima on the screen, before the plates were introduced is now occupied by the third bright fringe. Find the thickness of the plates (in $\mu$m), when the wavelength of light used is 4000\,\AA.
\qn{16} In a particular set up, interference is observed on a screen as shown. Wavelength of light emitted by source is 4000\,\AA. out of two coherent sources, one is obtained by reflection with a plane mirror of length 2 m. If $I_{\max}$ is the maximum intensity in interference pattern, find number of such points in the pattern.
\qn{17} Consider Bohr's theory for hydrogen atom. The magnitude of angular momentum, orbit radius and frequency of the electron in $n$th energy state in hydrogen atom are $L$, $r$ and $f$ respectively. Find out the value of `$x$', if the product $f\,r\,L$ directly proportional to $n^x$:
\qn{18} In a hydrogen like atom the energy required to excite the electron from 2nd to 3rd orbit is 47.2 eV. What is the atomic number of the atom?
\partsec{2}{CHEMISTRY}
\sect{SECTION-I (i)}
\qn{1} The products obtained when `Pt' reacts with aqua regia.
\begin{qoptions}
\opt{A}{$\mathrm{PtCl_4,\ NO_2}$}
\opt{B}{$\mathrm{H_2PtCl_4,\ N_2O}$}
\opt{C}{$\mathrm{H_2PtCl_6,\ NOCl}$}
\opt{D}{$\mathrm{PtCl_2,\ N_2O_3}$}
\end{qoptions}
\qn{2} Which of the following reaction is incorrect ?
\begin{qoptions}
\opt{A}{$\mathrm{P_4+8SOCl_2\to4PCl_3+4SO_2+2S_2Cl_2}$}
\opt{B}{$\mathrm{2I_{2(dry)}+3O_3\to I_4O_9}$}
\opt{C}{$\mathrm{3I^-_{aq}+O_3+2H^+\to I_3^-+O_2+H_2O}$}
\opt{D}{$\mathrm{AuCl_3+4Na_2S_2O_3\to}$\\ $\mathrm{Na_2\!\left(Au(S_2O_3)_2\right)+Na_2S_4O_6+3NaCl}$}
\end{qoptions}
\qn{3} Which of the following statement is incorrect ?
\begin{qoptions}
\opt{A}{Dilute acids decomposes $\mathrm{Na_2S_2O_3}$ to S and $\mathrm{SO_2}$}
\opt{B}{$\mathrm{H_2S_2O_4}$, $\mathrm{H_2S_2O_5}$ and $\mathrm{H_2S_2O_6}$ having S -- S linkages}
\opt{C}{Excess of $\mathrm{PCl_5}$ reacts with $\mathrm{H_2SO_4}$ to form a $\mathrm{SO_2Cl_2}$ and $\mathrm{POCl_3}$}
\opt{D}{Bleaching action of $\mathrm{SO_2}$ in the presence of moist is a oxidation process}
\end{qoptions}
\qn{4} Which reaction cannot be used for the preparation of the $\mathrm{ClO_2}$.
\begin{qoptions}
\opt{A}{$\mathrm{2AgClO_3+Cl_2\xrightarrow{90^\circ C}2AgCl+2ClO_2+O_2}$}
\opt{B}{$\mathrm{2NaClO_3+SO_2+H_2SO_4}\xrightarrow[\mathrm{NaCl}]{\mathrm{trace\ of}}$\\ $\mathrm{2ClO_2+2NaHSO_4}$}
\opt{C}{$\mathrm{3Cl_2+4HgO+\dfrac12O_2\xrightarrow{\Delta}2HgCl_2{\cdot}HgO+2ClO_2}$}
\opt{D}{$\mathrm{2NaClO_3+2H_2C_2O_4}\xrightarrow[90^\circ\mathrm{C}]{\mathrm{H_2O}}$\\ $\mathrm{2ClO_2+2CO_2+Na_2C_2O_4+2H_2O}$}
\end{qoptions}
\columnbreak
\sect{SECTION-I (ii)}
\qn{5} The correct statements among the following :
\begin{qoptions}
\opt{A}{$\mathrm{ClO_2}$ does not dimerize due to delocalisation of odd electron.}
\opt{B}{The ions of solid $\mathrm{I_2O_4}$ are $\mathrm{IO^+}$ and $\mathrm{IO_3^-}$ and ion of solid $\mathrm{I_4O_9}$ are $\mathrm{I^{3+}}$ and $\left(\mathrm{IO_3^-}\right)_3$.}
\opt{C}{Increasing oxidising power : $\mathrm{HClO_4}<\mathrm{HClO_3}<\mathrm{HClO_2}<\mathrm{HOCl}$}
\opt{D}{In $\mathrm{I_2Cl_2}$ structure bridging I -- Cl bonds are longer than terminal I -- Cl bonds.}
\end{qoptions}
\qn{6} Among the following the correct statements are :
\begin{qoptions}
\opt{A}{$\mathrm{ClF_3}$ is the most reactive inter halogen compound}
\opt{B}{Hydrolysis of $\mathrm{BrF_5}$ gives $\mathrm{HBrO_3}$ and HF}
\opt{C}{Decomposition of $\mathrm{Rb[ICl_2]}$ gives RbCl and ICl}
\opt{D}{In alkaline medium xenate disproportionates to perxenate and Xe.}
\end{qoptions}
\qfullbegin
\qn{7} The major product A and B in the following reaction sequence are :
\qfig{p1_chem_s2_q07_1}
\qfullsplit
\qfig[0.78]{p1_chem_s2_q07_A}
\qfullend
\qn{8} Which of the followings will give cyclic product upon heating ?
\begin{qoptions}
\opt{A}{$\mathrm{H_3C-CH_2-CH_2-CH_2-HC}\begin{array}{@{}l@{}}\scriptstyle\mathrm{COOH}\\[-2pt]\scriptstyle\mathrm{OH}\end{array}$}
\opt{B}{$\mathrm{H_3C-\underset{OH}{CH}-CH_2-CH_2-CH_2-COOH}$}
\opt{C}{$\mathrm{H_3C-CH_2-\underset{OH}{CH}-CH_2-CH_2-COOH}$}
\opt{D}{$\mathrm{H_3C-CH_2-CH_2-\underset{OH}{HC}-CH_2-COOH}$}
\end{qoptions}
\qn{9}
\qfig[0.8]{p1_chem_s2_q09_1}
\begin{qoptions}
\opt{A}{A will give +ve test with Tollen's reagent}
\opt{B}{A is a weak organic monobasic acid}
\opt{C}{A is a aldehyde}
\opt{D}{A on reaction with conc. $\mathrm{H_2SO_4}$ produces a neutral gas which is a very good reducing agent}
\end{qoptions}
\qn{10}
\qfig{p1_chem_s2_q10_1}
\begin{qoptions}
\opt{A}{Both A and B are aromatic in nature}
\opt{B}{Product A can show nucleophilic addition reaction}
\opt{C}{If potential nucleophile is not present then A will form dimer}
\opt{D}{B is a mixture of two structural isomers}
\end{qoptions}
\qn{11} $\mathrm{CH_3-CH_2-NO_2\,(X)\xrightarrow{Sn+HCl}K}$ (Organic compound), then which of the following is/are true for K ?
\begin{qoptions}
\opt{A}{K (form solid oxamide) with diethyl oxalate}
\opt{B}{K (form liquid oxamide) with diethyl oxalate}
\opt{C}{K will not give Liberman nitroso test}
\opt{D}{K will form blue colour complex with $\mathrm{Cu^{+2}}$ salt}
\end{qoptions}
\qfullbegin
\qn{12} Identify the correct statement(s) regarding the reductive ozonolysis of 2-pentene.
\begin{qoptions}
\opt{A}{A maximum of 10 ozonides can be formed}
\end{qoptions}
\qfig[0.8]{p1_chem_s2_q12_1}
\qfullend
\sect{SECTION-II}
\qn{13} How many 3d series elements (Sc to Zn) have positive standard potential for $\mathrm{M^{+2}/M}$.
\qn{14} The oxidation state of Mn in its highest fluoride.
\qn{15} The oxidation of the Cu in Schweizer's reagent which is obtained in the following reaction : $\mathrm{CuSO_4+NH_4OH_{(Excess)}\to}$ Schweizer's reagent $\mathrm{+H_2O}$
\qn{16} The enthalpy of hydration of $\mathrm{Cr^{+2}}$ ion is $-460$ k cal/mole (experimental value). The calculated value is $-435$ k cal/mole. What is the $\Delta_0$ for $\mathrm{[Cr(H_2O)_6]^{2+}}$ ion.
\qn{17} The number of metal - metal bonds in $\mathrm{Mn_2(CO)_{10}}$
\qn{18} The number of isomers possible for $\mathrm{[Pt\,(NH_3)_2(Py)_2Cl_2]^{2+}}$.
\partsec{3}{MATHEMATICS}
\sect{SECTION-I (i)}
\qn{1} A bag contains 8 red and 5 white balls. Three balls are drawn at random. The probability that one is red and two balls are white is
\begin{qoptions}
\opt{A}{$\dfrac{40}{143}$}
\opt{B}{$\dfrac{43}{286}$}
\opt{C}{$\dfrac{38}{143}$}
\opt{D}{$\dfrac{41}{286}$}
\end{qoptions}
\qn{2} $A$ and $B$ are independent events. The probability that both $A$ and $B$ occur is $\dfrac18$ and the probability that none of them occurs is $\dfrac38$. The probability of occurrence of event $A$ is $P_1$ or $P_2$. Then $P_1^2+P_2$ equals $(P_1<P_2)$
\begin{qoptions}
\opt{A}{$\dfrac7{16}$}
\opt{B}{$\dfrac9{16}$}
\opt{C}{$\dfrac58$}
\opt{D}{$\dfrac78$}
\end{qoptions}
\qn{3} $A$ is targeting $B$, $B$ and $C$ are targeting $A$. The probability of hitting the target by $A$, $B$ and $C$ are $\dfrac23,\dfrac12,\dfrac13$ respectively. If $A$ is hit, the probability that $B$ hits the target and $C$ does not is
\begin{qoptions}
\opt{A}{$\dfrac13$}
\opt{B}{$\dfrac12$}
\opt{C}{$\dfrac23$}
\opt{D}{$\dfrac56$}
\end{qoptions}
\qn{4} Three groups $A$, $B$ and $C$ are contesting for position on the board of directors of a company. The probabilities of their winning are 0.5, 0.3 and 0.2 respectively. If the group $A$ wins, then the probability of introducing a new product is 0.7 and the corresponding probabilities for groups $B$ and $C$ are 0.6 and 0.5 respectively. Find the probability that the new product will be introduced is
\begin{qoptions}
\opt{A}{0.6}
\opt{B}{0.61}
\opt{C}{0.62}
\opt{D}{0.63}
\end{qoptions}
\sect{SECTION-I (ii)}
\qn{5} In a test, a student either guesses or copies or knows the answer to a multiple choice question with four choices. The probability that he makes a guess is $\dfrac13$ and the probability that he copies the answer is $\dfrac16$. The probability that his answer is correct given that he copies it is $\dfrac18$. The probability that he knew the answer to the question given that he correctly answered it is $\dfrac xy$, where $x$ and $y$ are co-prime, then
\begin{qoptions}
\opt{A}{$y-x=6$}
\opt{B}{$x+y=54$}
\opt{C}{$2x-y$ is prime}
\opt{D}{The units digit of $x$ is divisible by 4}
\end{qoptions}
\qn{6} If on an average out of 10 ships, one is drowned, then the probability that out of 5 ships, at least 4 reach safely is $2^a\times3^b\times5^c\times7^d$, then (where $a,b,c,d$ are integers)
\begin{qoptions}
\opt{A}{$a+b+c+d=0$}
\opt{B}{$b+d$ is prime}
\opt{C}{$a+b$ is a perfect square}
\opt{D}{$b+3d+c=5$}
\end{qoptions}
\qn{7} A person throws two dice, one a common cube with faces: 1, 2, 3, 4, 5, 6 and other a regular tetrahedron with faces 1, 2, 3, 4. Then the probability that the sum of numbers appearing on the dice is 6 is $\dfrac ab$ where $a$ and $b$ are co-prime integers, then
\begin{qoptions}
\opt{A}{$0.1<\dfrac ab<0.2$}
\opt{B}{$a$ is prime and $b$ is composite}
\opt{C}{$a+b$ divides 72}
\opt{D}{$2a-b$ is negative}
\end{qoptions}
\qn{8} The probability that in a year of 22$^{nd}$ century choosen at random, their will be 53 sundays, is $\dfrac pq$ where $p$ and $q$ are co-prime, then
\begin{qoptions}
\opt{A}{$0.2<\dfrac pq<0.3$}
\opt{B}{$p^2+q=53$}
\opt{C}{$p+q$ is divisible by 11}
\opt{D}{$q$ is divisible by 7}
\end{qoptions}
\qn{9} Two numbers $b$ and $c$ are chosen at random with replacement from the numbers 1, 2, 3, 4, 5, 6, 7, 8, 9. The probability that $x^2+bx+c>0\ \forall x\in R$ is $\dfrac mn$ where $m$ and $n$ are co-prime integers. Then
\begin{qoptions}
\opt{A}{Number of divisors of $m$ is 6}
\opt{B}{$m^2+n$ exceeds 1000}
\opt{C}{Sum of digits of $m$+sum of digits of $n$ is 13}
\opt{D}{$m$ is divisible by 8}
\end{qoptions}
\qn{10} Three digit numbers are formed using the digits 0, 1, 2, 3, 4, 5 without repetition of digits. If a number is choosen at random, then the probability that the digits either increase or decreases is $\dfrac\alpha\beta$ where $\alpha$ and $\beta$ are co-prime integers, then
\begin{qoptions}
\opt{A}{$\alpha$ is prime and $\beta$ is composite}
\opt{B}{$\alpha^2<\beta$}
\opt{C}{$\alpha+2\beta$ has 3 factors}
\opt{D}{$\beta-2\alpha$ is a perfect square}
\end{qoptions}
\qn{11} Let $P_n$ denote the probability of getting $n$ heads when a fair coin is tossed $m$ times $(n\le m)$. If $P_4,P_5,P_6$ are in AP, then $m$ can be
\begin{qoptions}
\opt{A}{7}
\opt{B}{10}
\opt{C}{14}
\opt{D}{16}
\end{qoptions}
\qn{12} Three critics review a book. Odds in favour of the book are 5 : 2, 4 : 3 and 3 : 4 respectively for the 3 critics. The probability that majority are in the favour of the book is $\dfrac uv$, where $u$ \& $v$ are co-prime integers, then
\begin{qoptions}
\opt{A}{$u$ has 4 divisors}
\opt{B}{The sum of digits of $u$ is greater than the sum of digits of $v$}
\opt{C}{One of the digits of $u$ is a square of one of the digits of $v$}
\opt{D}{One of the digits of $v$ is a square of one of the digits of $u$}
\end{qoptions}
\sect{SECTION-II}
\qn{13} Let $w$ be the imaginary cube root of unity $(w\ne1)$. A fair dice is thrown 3 times. If $r_1,r_2$ and $r_3$ are the numbers obtained on the dice, then the probability that $w^{r_1}+w^{r_2}+w^{r_3}=0$ is $\dfrac hk$, where $h$ and $k$ are co-prime integers, then $h^2+k^3$ equals \underline{\hspace{1.5cm}}
\qn{14} An unbiased dice is tossed until a number greater than 4 appears. The probability that an even number of tosses is needed is $\dfrac\alpha\beta$, where $\alpha$ and $\beta$ are co-prime integers, then $\alpha^\beta+\beta^\alpha$ equals \underline{\hspace{1.5cm}}
\qn{15} A box contains 10 mangoes out of which 4 are rotten. Two mangoes are taken out together. If one of them is found to be good, then the probability that the other mango is also good is $\dfrac xy$, where $x$ and $y$ are co-prime integers, then $y-2x$ equals \underline{\hspace{1.5cm}}
\qn{16} A bag contains 7 balls out of which some are white. The probability that the bag contains exactly $x$ white balls is proportional to $x^2\ (1\le x\le7)$. A ball is drawn from the bag and found to be white. The probability that the bag contains 2 white balls is $\dfrac ab$ where $a$ and $b$ are co-prime integers, then $2a^2+b$ equals \underline{\hspace{1.5cm}}
\qn{17} Of the three independent events $E_1,E_2$ and $E_3$, the probability that only $E_1$ occurs is $\alpha$, only $E_2$ occurs is $\beta$ and only $E_3$ occurs is $\gamma$. Let the probability $P$ that none of the events $E_1,E_2$ or $E_3$ occurs satisfy $(\alpha-2\beta)P=\alpha\beta$ and $(\beta-3\gamma)P=2\beta\gamma$. All the given probabilities are assumed to lie in the interval $(0,1)$. Then $\dfrac{\text{Probability of occurence of }E_1}{\text{Probability of occurence of }E_3}$ equals \underline{\hspace{1.5cm}}
\qn{18} If $P(B)=\dfrac34$, $P\left(A\cap B\cap\overline C\right)=\dfrac13$ and $p\left(\overline A\cap B\cap\overline C\right)=\dfrac13$, then $P(B\cap C)$ equals $\dfrac1h$, then $h$ equals \underline{\hspace{1.5cm}}
\end{multicols}
\end{document}
